% Source: book pp. 298--307 (PDF 312--321). No solutions.
% numbered displays (Axler's 8.5, 8.7, ...): step the item counter and label
% it just before the display, then \tag*{\thetheorem} inside the display
\DeclareRobustCommand\axeqnum[1]{\refstepcounter{theorem}\label{ax:#1}}
% the italic comment that follows some exercises

\section{Generalized Eigenvectors and Nilpotent Operators}

\subsection{Null Spaces of Powers of an Operator}

We begin this chapter with a study of null spaces of powers of an operator.

\begin{result}[Sequence of increasing null spaces]\label{ax:8.1}
Suppose $T\in\Lin(V)$. Then
\[ \{0\} = \nul T^0 \subseteq \nul T^1 \subseteq \cdots \subseteq \nul T^k
   \subseteq \nul T^{k+1} \subseteq \cdots. \]
\end{result}

\begin{proof}
Suppose $k$ is a nonnegative integer and $v\in\nul T^k$. Then $T^kv=0$, which
implies that $T^{k+1}v=T(T^kv)=T(0)=0$. Thus $v\in\nul T^{k+1}$. Hence
$\nul T^k\subseteq\nul T^{k+1}$, as desired.
\end{proof}

\marginnote{For similar results about decreasing sequences of ranges, see
Exercises 6, 7, and 8.}%
The following result states that if two consecutive terms in the sequence of
subspaces above are equal, then all later terms in the sequence are equal.

\begin{result}[Equality in the sequence of null spaces]\label{ax:8.2}
Suppose $T\in\Lin(V)$ and $m$ is a nonnegative integer such that
\[ \nul T^m = \nul T^{m+1}. \]
Then
\[ \nul T^m = \nul T^{m+1} = \nul T^{m+2} = \nul T^{m+3} = \cdots. \]
\end{result}

\begin{proof}
Let $k$ be a positive integer. We want to prove that
\[ \nul T^{m+k} = \nul T^{m+k+1}. \]
We already know from \ref{ax:8.1} that $\nul T^{m+k}\subseteq\nul T^{m+k+1}$.

To prove the inclusion in the other direction, suppose $v\in\nul T^{m+k+1}$.
Then
\[ T^{m+1}(T^kv) = T^{m+k+1}v = 0. \]
Hence
\[ T^kv\in\nul T^{m+1} = \nul T^m. \]
Thus $T^{m+k}v=T^m(T^kv)=0$, which means that $v\in\nul T^{m+k}$. This implies
that $\nul T^{m+k+1}\subseteq\nul T^{m+k}$, completing the proof.
\end{proof}

The result above raises the question of whether there exists a nonnegative
integer $m$ such that $\nul T^m=\nul T^{m+1}$. The next result shows that this
equality holds at least when $m$ equals the dimension of the vector space on
which $T$ operates.

\begin{result}[Null spaces stop growing]\label{ax:8.3}
Suppose $T\in\Lin(V)$. Then
\[ \nul T^{\dim V} = \nul T^{\dim V+1} = \nul T^{\dim V+2} = \cdots. \]
\end{result}

\begin{proof}\emergencystretch=1.5em
We only need to prove that $\nul T^{\dim V}=\nul T^{\dim V+1}$ (by
\ref{ax:8.2}). Suppose this is not true. Then, by \ref{ax:8.1} and
\ref{ax:8.2}, we have
\[ \{0\} = \nul T^0 \subsetneq \nul T^1 \subsetneq \cdots \subsetneq
   \nul T^{\dim V} \subsetneq \nul T^{\dim V+1}, \]
where the symbol $\subsetneq$ means ``contained in but not equal to''. At each
of the strict inclusions in the chain above, the dimension increases by at
least $1$. Thus $\dim\nul T^{\dim V+1}\ge\dim V+1$, a contradiction because a
subspace of $V$ cannot have a larger dimension than $\dim V$.
\end{proof}

It is not true that $V=\nul T\oplus\range T$ for every $T\in\Lin(V)$. However,
the next result can be a useful substitute.

\begin{result}[$V$ is the direct sum of $\nul T^{\dim V}$ and
$\range T^{\dim V}$]\label{ax:8.4}
Suppose $T\in\Lin(V)$. Then
\[ V = \nul T^{\dim V}\oplus\range T^{\dim V}. \]
\end{result}

\begin{proof}
Let $n=\dim V$. First we show that\axeqnum{8.5}%
\[ (\nul T^n)\cap(\range T^n) = \{0\}. \tag*{\thetheorem} \]
Suppose $v\in(\nul T^n)\cap(\range T^n)$. Then $T^nv=0$, and there exists
$u\in V$ such that $v=T^nu$. Applying $T^n$ to both sides of the last equation
shows that $T^nv=T^{2n}u$. Hence $T^{2n}u=0$, which implies that $T^nu=0$ (by
\ref{ax:8.3}). Thus $v=T^nu=0$, completing the proof of \ref{ax:8.5}.

Now \ref{ax:8.5} implies that $\nul T^n+\range T^n$ is a direct sum (by
\ref{ax:1.46}). Also,
\[ \dim(\nul T^n\oplus\range T^n) = \dim\nul T^n+\dim\range T^n = \dim V, \]
where the first equality above comes from \ref{ax:3.94} and the second equality
comes from the fundamental theorem of linear maps (\ref{ax:3.21}). The equation
above implies that $\nul T^n\oplus\range T^n=V$ (see \ref{ax:2.39}), as
desired.
\end{proof}

For an improvement of the result above, see Exercise 19.

\begin{example}[$\F^3=\nul T^3\oplus\range T^3$ for $T\in\Lin(\F^3)$]\label{ax:8.6}
Suppose $T\in\Lin(\F^3)$ is defined by
\[ T(z_1,z_2,z_3) = (4z_2,0,5z_3). \]
Then $\nul T=\{(z_1,0,0) : z_1\in\F\}$ and
$\range T=\{(z_1,0,z_3) : z_1,z_3\in\F\}$. Thus $\nul T\cap\range T\ne\{0\}$.
Hence $\nul T+\range T$ is not a direct sum. Also note that
$\nul T+\range T\ne\F^3$. However, we have $T^3(z_1,z_2,z_3)=(0,0,125z_3)$.
Thus we see that
\[ \nul T^3 = \{(z_1,z_2,0) : z_1,z_2\in\F\} \quad\text{and}\quad
   \range T^3 = \{(0,0,z_3) : z_3\in\F\}. \]
Hence $\F^3=\nul T^3\oplus\range T^3$, as expected by \ref{ax:8.4}.
\end{example}

\subsection{Generalized Eigenvectors}

Some operators do not have enough eigenvectors to lead to good descriptions of
their behavior. Thus in this subsection we introduce the concept of generalized
eigenvectors, which will play a major role in our description of the structure
of an operator.

To understand why we need more than eigenvectors, let's examine the question of
describing an operator by decomposing its domain into invariant subspaces. Fix
$T\in\Lin(V)$. We seek to describe $T$ by finding a ``nice'' direct sum
decomposition
\[ V = V_1\oplus\cdots\oplus V_n, \]
where each $V_k$ is a subspace of $V$ invariant under $T$. The simplest
possible nonzero invariant subspaces are one-dimensional. A decomposition as
above in which each $V_k$ is a one-dimensional subspace of $V$ invariant under
$T$ is possible if and only if $V$ has a basis consisting of eigenvectors of
$T$ (see \ref{ax:5.55}). This happens if and only if $V$ has an eigenspace
decomposition\axeqnum{8.7}%
\[ V = E(\lambda_1,T)\oplus\cdots\oplus E(\lambda_m,T), \tag*{\thetheorem} \]
where $\lambda_1,\dots,\lambda_m$ are the distinct eigenvalues of $T$ (see
\ref{ax:5.55}).

The spectral theorem in the previous chapter shows that if $V$ is an inner
product space, then a decomposition of the form \ref{ax:8.7} holds for every
self-adjoint operator if $\F=\R$ and for every normal operator if $\F=\C$
because operators of those types have enough eigenvectors to form a basis of
$V$ (see \ref{ax:7.29} and \ref{ax:7.31}).

However, a decomposition of the form \ref{ax:8.7} may not hold for more general
operators, even on a complex vector space. An example was given by the operator
in \ref{ax:5.57}, which does not have enough eigenvectors for \ref{ax:8.7} to
hold. Generalized eigenvectors and generalized eigenspaces, which we now
introduce, will remedy this situation.

\begin{definition}[Generalized eigenvector]\label{ax:8.8}
Suppose $T\in\Lin(V)$ and $\lambda$ is an eigenvalue of $T$. A vector $v\in V$
is called a \emph{generalized eigenvector} of $T$ corresponding to $\lambda$ if
$v\ne0$ and
\[ (T-\lambda I)^kv = 0 \]
for some positive integer $k$.
\end{definition}

\marginnote{Generalized eigenvalues are not defined because doing so would not
lead to anything new. Reason: if $(T-\lambda I)^k$ is not injective for some
positive integer $k$, then $T-\lambda I$ is not injective, and hence $\lambda$
is an eigenvalue of $T$.}%
A nonzero vector $v\in V$ is a generalized eigenvector of $T$ corresponding to
$\lambda$ if and only if
\[ (T-\lambda I)^{\dim V}v = 0, \]
as follows from applying \ref{ax:8.1} and \ref{ax:8.3} to the operator
$T-\lambda I$.

As we know, an operator on a complex vector space may not have enough
eigenvectors to form a basis of the domain. The next result shows that on a
complex vector space there are enough generalized eigenvectors to do this.

\begin{result}[A basis of generalized eigenvectors]\label{ax:8.9}
Suppose $\F=\C$ and $T\in\Lin(V)$. Then there is a basis of $V$ consisting of
generalized eigenvectors of $T$.
\end{result}

\begin{proof}
Let $n=\dim V$. We will use induction on $n$. To get started, note that the
desired result holds if $n=1$ because then every nonzero vector in $V$ is an
eigenvector of $T$.

\marginnote{This step is where we use the hypothesis that $\F=\C$, because if
$\F=\R$ then $T$ may not have any eigenvalues.}%
Now suppose $n>1$ and the desired result holds for all smaller values of
$\dim V$. Let $\lambda$ be an eigenvalue of $T$. Applying \ref{ax:8.4} to
$T-\lambda I$ shows that
\[ V = \nul(T-\lambda I)^n\oplus\range(T-\lambda I)^n. \]

If $\nul(T-\lambda I)^n=V$, then every nonzero vector in $V$ is a generalized
eigenvector of $T$, and thus in this case there is a basis of $V$ consisting
of generalized eigenvectors of $T$. Hence we can assume that
$\nul(T-\lambda I)^n\ne V$, which implies that
$\range(T-\lambda I)^n\ne\{0\}$.

Also, $\nul(T-\lambda I)^n\ne\{0\}$, because $\lambda$ is an eigenvalue of $T$.
Thus we have
\[ 0 < \dim\range(T-\lambda I)^n < n. \]
Furthermore, $\range(T-\lambda I)^n$ is invariant under $T$
$\bigl[$by \ref{ax:5.18} with $p(z)=(z-\lambda)^n\bigr]$. Let
$S\in\Lin\bigl(\range(T-\lambda I)^n\bigr)$ equal $T$ restricted to
$\range(T-\lambda I)^n$. Our induction hypothesis applied to the operator $S$
implies that there is a basis of $\range(T-\lambda I)^n$ consisting of
generalized eigenvectors of $S$, which of course are generalized eigenvectors
of $T$. Adjoining that basis of $\range(T-\lambda I)^n$ to a basis of
$\nul(T-\lambda I)^n$ gives a basis of $V$ consisting of generalized
eigenvectors of $T$.
\end{proof}

If $\F=\R$ and $\dim V>1$, then some operators on $V$ have the property that
there exists a basis of $V$ consisting of generalized eigenvectors of the
operator, and (unlike what happens when $\F=\C$) other operators do not have
this property. See Exercise 11 for a necessary and sufficient condition that
determines whether an operator has this property.

\begin{example}[Generalized eigenvectors of an operator on $\C^3$]\label{ax:8.10}
Define $T\in\Lin(\C^3)$ by
\[ T(z_1,z_2,z_3) = (4z_2,0,5z_3) \]
for each $(z_1,z_2,z_3)\in\C^3$. A routine use of the definition of eigenvalue
shows that the eigenvalues of $T$ are $0$ and $5$. Furthermore, the
eigenvectors corresponding to the eigenvalue $0$ are the nonzero vectors of the
form $(z_1,0,0)$, and the eigenvectors corresponding to the eigenvalue $5$ are
the nonzero vectors of the form $(0,0,z_3)$. Hence this operator does not have
enough eigenvectors to span its domain $\C^3$.

We compute that $T^3(z_1,z_2,z_3)=(0,0,125z_3)$. Thus \ref{ax:8.1} and
\ref{ax:8.3} imply that the generalized eigenvectors of $T$ corresponding to the
eigenvalue $0$ are the nonzero vectors of the form $(z_1,z_2,0)$.

We also have $(T-5I)^3(z_1,z_2,z_3)=(-125z_1+300z_2,-125z_2,0)$. Thus the
generalized eigenvectors of $T$ corresponding to the eigenvalue $5$ are the
nonzero vectors of the form $(0,0,z_3)$.

The paragraphs above show that each of the standard basis vectors of $\C^3$ is a
generalized eigenvector of $T$. Thus $\C^3$ indeed has a basis consisting of
generalized eigenvectors of $T$, as promised by \ref{ax:8.9}.
\end{example}

If $v$ is an eigenvector of $T\in\Lin(V)$, then the corresponding eigenvalue
$\lambda$ is uniquely determined by the equation $Tv=\lambda v$, which can be
satisfied by only one $\lambda\in\F$ (because $v\ne0$). However, if $v$ is a
generalized eigenvector of $T$, then it is not obvious that the equation
$(T-\lambda I)^{\dim V}v=0$ can be satisfied by only one $\lambda\in\F$.
Fortunately, the next result tells us that all is well on this issue.

\begin{result}[Generalized eigenvector corresponds to a unique
eigenvalue]\label{ax:8.11}
Suppose $T\in\Lin(V)$. Then each generalized eigenvector of $T$ corresponds to
only one eigenvalue of $T$.
\end{result}

\begin{proof}
Suppose $v\in V$ is a generalized eigenvector of $T$ corresponding to
eigenvalues $\alpha$ and $\lambda$ of $T$. Let $m$ be the smallest positive
integer such that $(T-\alpha I)^mv=0$. Let $n=\dim V$. Then
\begin{align*}
0 &= (T-\lambda I)^nv\\
  &= \bigl((T-\alpha I)+(\alpha-\lambda)I\bigr)^nv\\
  &= \sum_{k=0}^{n} b_k(\alpha-\lambda)^{n-k}(T-\alpha I)^kv,
\end{align*}
where $b_0=1$ and the values of the other binomial coefficients $b_k$ do not
matter. Apply the operator $(T-\alpha I)^{m-1}$ to both sides of the equation
above, getting
\[ 0 = (\alpha-\lambda)^n(T-\alpha I)^{m-1}v. \]
Because $(T-\alpha I)^{m-1}v\ne0$, the equation above implies that
$\alpha=\lambda$, as desired.
\end{proof}

We saw earlier (\ref{ax:5.11}) that eigenvectors corresponding to distinct
eigenvalues are linearly independent. Now we prove a similar result for
generalized eigenvectors, with a proof that roughly follows the pattern of the
proof of that earlier result.

\begin{result}[Linearly independent generalized eigenvectors]\label{ax:8.12}
Suppose that $T\in\Lin(V)$. Then every list of generalized eigenvectors of $T$
corresponding to distinct eigenvalues of $T$ is linearly independent.
\end{result}

\begin{proof}
Suppose the desired result is false. Then there exists a smallest positive
integer $m$ such that there exists a linearly dependent list
$v_1,\dots,v_m$ of generalized eigenvectors of $T$ corresponding to distinct
eigenvalues $\lambda_1,\dots,\lambda_m$ of $T$ (note that $m\ge2$ because a
generalized eigenvector is, by definition, nonzero). Thus there exist
$a_1,\dots,a_m\in\F$, none of which are $0$ (because of the minimality of $m$),
such that
\[ a_1v_1+\cdots+a_mv_m = 0. \]

Let $n=\dim V$. Apply $(T-\lambda_mI)^n$ to both sides of the equation above,
getting\axeqnum{8.13}%
\[ a_1(T-\lambda_mI)^nv_1+\cdots+a_{m-1}(T-\lambda_mI)^nv_{m-1} = 0.
   \tag*{\thetheorem} \]
Suppose $k\in\{1,\dots,m-1\}$. Then
\[ (T-\lambda_mI)^nv_k \ne 0 \]
because otherwise $v_k$ would be a generalized eigenvector of $T$ corresponding
to the distinct eigenvalues $\lambda_k$ and $\lambda_m$, which would contradict
\ref{ax:8.11}. However,
\[ (T-\lambda_kI)^n\bigl((T-\lambda_mI)^nv_k\bigr)
   = (T-\lambda_mI)^n\bigl((T-\lambda_kI)^nv_k\bigr) = 0. \]

Thus the last two displayed equations show that $(T-\lambda_mI)^nv_k$ is a
generalized eigenvector of $T$ corresponding to the eigenvalue $\lambda_k$.
Hence
\[ (T-\lambda_mI)^nv_1,\dots,(T-\lambda_mI)^nv_{m-1} \]
is a linearly dependent list (by \ref{ax:8.13}) of $m-1$ generalized
eigenvectors corresponding to distinct eigenvalues, contradicting the
minimality of $m$. This contradiction completes the proof.
\end{proof}

\subsection{Nilpotent Operators}

\begin{definition}[Nilpotent]\label{ax:8.14}
An operator is called \emph{nilpotent} if some power of it equals $0$.
\end{definition}

Thus an operator $T\in\Lin(V)$ is nilpotent if and only if every nonzero vector
in $V$ is a generalized eigenvector of $T$ corresponding to the eigenvalue $0$.

\begin{example}[Nilpotent operators]\label{ax:8.15}\leavevmode
\begin{parts}
\item The operator $T\in\Lin(\F^4)$ defined by
  \[ T(z_1,z_2,z_3,z_4) = (0,0,z_1,z_2) \]
  is nilpotent because $T^2=0$.
\item The operator on $\F^3$ whose matrix (with respect to the standard basis)
  is
  \[ \begin{pmatrix} -3 & 9 & 0\\ -7 & 9 & 6\\ 4 & 0 & -6 \end{pmatrix} \]
  is nilpotent, as can be shown by cubing the matrix above to get the zero
  matrix.
\item The operator of differentiation on $\Poly_m(\R)$ is nilpotent because the
  $(m+1)^{\text{th}}$ derivative of every polynomial of degree at most $m$
  equals $0$. Note that on this space of dimension $m+1$, we need to raise the
  nilpotent operator to the power $m+1$ to get the $0$ operator.
\end{parts}
\end{example}

\marginnote{The Latin word \textbf{nil} means nothing or zero; the Latin word
\textbf{potens} means having power. Thus \textbf{nilpotent} literally means
having a power that is zero.}%
The next result shows that when raising a nilpotent operator to a power, we
never need to use a power higher than the dimension of the space. For a
slightly stronger result, see Exercise 18.

\begin{result}[Nilpotent operator raised to dimension of domain is
$0$]\label{ax:8.16}
Suppose $T\in\Lin(V)$ is nilpotent. Then $T^{\dim V}=0$.
\end{result}

\begin{proof}
Because $T$ is nilpotent, there exists a positive integer $k$ such that
$T^k=0$. Thus $\nul T^k=V$. Now \ref{ax:8.1} and \ref{ax:8.3} imply that
$\nul T^{\dim V}=V$. Thus $T^{\dim V}=0$.
\end{proof}

\begin{result}[Eigenvalues of nilpotent operator]\label{ax:8.17}
Suppose $T\in\Lin(V)$.
\begin{parts}
\item If $T$ is nilpotent, then $0$ is an eigenvalue of $T$ and $T$ has no
  other eigenvalues.
\item If $\F=\C$ and $0$ is the only eigenvalue of $T$, then $T$ is nilpotent.
\end{parts}
\end{result}

\begin{proof}\leavevmode
\begin{parts}
\item To prove (a), suppose $T$ is nilpotent. Hence there is a positive integer
  $m$ such that $T^m=0$. This implies that $T$ is not injective. Thus $0$ is an
  eigenvalue of $T$.

  To show that $T$ has no other eigenvalues, suppose $\lambda$ is an eigenvalue
  of $T$. Then there exists a nonzero vector $v\in V$ such that
  \[ \lambda v = Tv. \]
  Repeatedly applying $T$ to both sides of this equation shows that
  \[ \lambda^mv = T^mv = 0. \]
  Thus $\lambda=0$, as desired.
\item Suppose $\F=\C$ and $0$ is the only eigenvalue of $T$. By
  \ref{ax:5.27}(b), the minimal polynomial of $T$ equals $z^m$ for some
  positive integer $m$. Thus $T^m=0$. Hence $T$ is nilpotent.
\end{parts}
\end{proof}

Exercise 23 shows that the hypothesis that $\F=\C$ cannot be deleted in (b) of
the result above.

Given an operator on $V$, we want to find a basis of $V$ such that the matrix
of the operator with respect to this basis is as simple as possible, meaning
that the matrix contains many $0$'s. The next result shows that if $T$ is
nilpotent, then we can choose a basis of $V$ such that the matrix of $T$ with
respect to this basis has more than half of its entries equal to $0$. Later in
this chapter we will do even better.

\begin{result}[Minimal polynomial and upper-triangular matrix of nilpotent
operator]\label{ax:8.18}
Suppose $T\in\Lin(V)$. Then the following are equivalent.
\begin{parts}
\item $T$ is nilpotent.
\item The minimal polynomial of $T$ is $z^m$ for some positive integer $m$.
\item There is a basis of $V$ with respect to which the matrix of $T$ has the
  form
  \[ \begin{pmatrix} 0 & & *\\ & \ddots & \\ 0 & & 0 \end{pmatrix}, \]
  where all entries on and below the diagonal equal $0$.
\end{parts}
\end{result}

\begin{proof}
Suppose (a) holds, so $T$ is nilpotent. Thus there exists a positive integer
$n$ such that $T^n=0$. Now \ref{ax:5.29} implies that $z^n$ is a polynomial
multiple of the minimal polynomial of $T$. Thus the minimal polynomial of $T$
is $z^m$ for some positive integer $m$, proving that (a) implies (b).

Now suppose (b) holds, so the minimal polynomial of $T$ is $z^m$ for some
positive integer $m$. This implies, by \ref{ax:5.27}(a), that $0$ (which is the
only zero of $z^m$) is the only eigenvalue of $T$. This further implies, by
\ref{ax:5.44}, that there is a basis of $V$ with respect to which the matrix of
$T$ is upper triangular. This also implies, by \ref{ax:5.41}, that all entries
on the diagonal of this matrix are $0$, proving that (b) implies (c).

Now suppose (c) holds. Then \ref{ax:5.40} implies that $T^{\dim V}=0$. Thus $T$
is nilpotent, proving that (c) implies (a).
\end{proof}

% ------------------------------------------------------------------ exercises
\exercisesheading{8A}

\begin{exercise}
Suppose $T\in\Lin(V)$. Prove that if $\dim\nul T^4=8$ and $\dim\nul T^6=9$,
then $\dim\nul T^m=9$ for all integers $m\ge5$.
\end{exercise}

\begin{exercise}
Suppose $T\in\Lin(V)$, $m$ is a positive integer, $v\in V$, and
$T^{m-1}v\ne0$ but $T^mv=0$. Prove that $v,Tv,T^2v,\dots,T^{m-1}v$ is linearly
independent.
\exnote{The result in this exercise is used in the proof of \ref{ax:8.45}.}
\end{exercise}

\begin{exercise}
Suppose $T\in\Lin(V)$. Prove that
\[ V = \nul T\oplus\range T \iff \nul T^2 = \nul T. \]
\end{exercise}

\begin{exercise}
Suppose $T\in\Lin(V)$, $\lambda\in\F$, and $m$ is a positive integer such that
the minimal polynomial of $T$ is a polynomial multiple of $(z-\lambda)^m$.
Prove that
\[ \dim\nul(T-\lambda I)^m \ge m. \]
\end{exercise}

\begin{exercise}
Suppose $T\in\Lin(V)$ and $m$ is a positive integer. Prove that
\[ \dim\nul T^m \le m\dim\nul T. \]
\exnote{Hint: Exercise 21 in Section 3B may be useful.}
\end{exercise}

\begin{exercise}
Suppose $T\in\Lin(V)$. Show that
\[ V = \range T^0 \supseteq \range T^1 \supseteq \cdots \supseteq \range T^k
   \supseteq \range T^{k+1} \supseteq \cdots. \]
\end{exercise}

\begin{exercise}
Suppose $T\in\Lin(V)$ and $m$ is a nonnegative integer such that
\[ \range T^m = \range T^{m+1}. \]
Prove that $\range T^k=\range T^m$ for all $k>m$.
\end{exercise}

\begin{exercise}
Suppose $T\in\Lin(V)$. Prove that
\[ \range T^{\dim V} = \range T^{\dim V+1} = \range T^{\dim V+2} = \cdots. \]
\end{exercise}

\begin{exercise}
Suppose $T\in\Lin(V)$ and $m$ is a nonnegative integer. Prove that
\[ \nul T^m = \nul T^{m+1} \iff \range T^m = \range T^{m+1}. \]
\end{exercise}

\begin{exercise}
Define $T\in\Lin(\C^2)$ by $T(w,z)=(z,0)$. Find all generalized eigenvectors
of $T$.
\end{exercise}

\begin{exercise}
Suppose that $T\in\Lin(V)$. Prove that there is a basis of $V$ consisting of
generalized eigenvectors of $T$ if and only if the minimal polynomial of $T$
equals $(z-\lambda_1)\cdots(z-\lambda_m)$ for some
$\lambda_1,\dots,\lambda_m\in\F$.
\exnote{Assume $\F=\R$ because the case $\F=\C$ follows from
\ref{ax:5.27}\textup{(b)} and \ref{ax:8.9}.}
\exnote{This exercise states that the condition for there to be a basis of $V$
consisting of generalized eigenvectors of $T$ is the same as the condition for
there to be a basis with respect to which $T$ has an upper-triangular matrix
(see \ref{ax:5.44}).}
\exnote{\textbf{Caution:} If $T$ has an upper-triangular matrix with respect to
a basis $v_1,\dots,v_n$ of $V$, then $v_1$ is an eigenvector of $T$ but it is
not necessarily true that $v_2,\dots,v_n$ are generalized eigenvectors of $T$.}
\end{exercise}

\begin{exercise}
Suppose $T\in\Lin(V)$ is such that every nonzero vector in $V$ is a
generalized eigenvector of $T$. Prove that there exists $\lambda\in\F$ such
that $T-\lambda I$ is nilpotent.
\end{exercise}

\begin{exercise}
Suppose $S,T\in\Lin(V)$ and $ST$ is nilpotent. Prove that $TS$ is nilpotent.
\end{exercise}

\begin{exercise}
Suppose $T\in\Lin(V)$ is nilpotent and $T\ne0$. Prove $T$ is not
diagonalizable.
\end{exercise}

\begin{exercise}
Suppose $\F=\C$ and $T\in\Lin(V)$. Prove that $T$ is diagonalizable if and only
if every generalized eigenvector of $T$ is an eigenvector of $T$.
\exnote{For $\F=\C$, this exercise adds another equivalence to the list of
conditions for diagonalizability in \ref{ax:5.55}.}
\end{exercise}

\begin{exercise}\leavevmode
\begin{parts}
\item Give an example of nilpotent operators $S,T$ on the same vector space
  such that neither $S+T$ nor $ST$ is nilpotent.
\item Suppose $S,T\in\Lin(V)$ are nilpotent and $ST=TS$. Prove that $S+T$ and
  $ST$ are nilpotent.
\end{parts}
\end{exercise}

\begin{exercise}
Suppose $T\in\Lin(V)$ is nilpotent and $m$ is a positive integer such that
$T^m=0$.
\begin{parts}
\item Prove that $I-T$ is invertible and that
  $(I-T)^{-1}=I+T+\cdots+T^{m-1}$.
\item Explain how you would guess the formula above.
\end{parts}
\end{exercise}

\begin{exercise}
Suppose $T\in\Lin(V)$ is nilpotent. Prove that $T^{1+\dim\range T}=0$.
\exnote{If $\dim\range T<\dim V-1$, then this exercise improves
\ref{ax:8.16}.}
\end{exercise}

\begin{exercise}
Suppose $T\in\Lin(V)$ is not nilpotent. Show that
\[ V = \nul T^{\dim V-1}\oplus\range T^{\dim V-1}. \]
\exnote{For operators that are not nilpotent, this exercise improves
\ref{ax:8.4}.}
\end{exercise}

\begin{exercise}
Suppose $V$ is an inner product space and $T\in\Lin(V)$ is normal and
nilpotent. Prove that $T=0$.
\end{exercise}

\begin{exercise}
Suppose $T\in\Lin(V)$ is such that $\nul T^{\dim V-1}\ne\nul T^{\dim V}$.
Prove that $T$ is nilpotent and that $\dim\nul T^k=k$ for every integer $k$
with $0\le k\le\dim V$.
\end{exercise}

\begin{exercise}
Suppose $T\in\Lin(\C^5)$ is such that $\range T^4\ne\range T^5$. Prove that $T$
is nilpotent.
\end{exercise}

\begin{exercise}
Give an example of an operator $T$ on a finite-dimensional real vector space
such that $0$ is the only eigenvalue of $T$ but $T$ is not nilpotent.
\exnote{This exercise shows that the implication \textup{(b)}
$\Longrightarrow$ \textup{(a)} in \ref{ax:8.17} does not hold without the
hypothesis that $\F=\C$.}
\end{exercise}

\begin{exercise}
For each item in Example \ref{ax:8.15}, find a basis of the domain vector space
such that the matrix of the nilpotent operator with respect to that basis has
the upper-triangular form promised by \ref{ax:8.18}(c).
\end{exercise}

\begin{exercise}
Suppose that $V$ is an inner product space and $T\in\Lin(V)$ is nilpotent.
Show that there is an orthonormal basis of $V$ with respect to which the matrix
of $T$ has the upper-triangular form promised by \ref{ax:8.18}(c).
\end{exercise}
